% 天蝎座（综述）
% license CCBYSA3
% type Wiki

本文根据 CC-BY-SA 协议转载翻译自维基百科\href{https://en.wikipedia.org/wiki/Scorpius}{相关文章}。

天蝎座（Scorpius）是一个黄道星座，位于南天半球，处在银河系中心附近，其西侧为天秤座，东侧为人马座。天蝎座是一组极为古老的星座，其被认识的历史早于希腊文化$^{[2]}$；它也是公元$^{[2]}$世纪希腊天文学家托勒密所确认的 48 个星座之一。
\subsection{显著特征}  
\subsubsection{恒星}
\begin{figure}[ht]
\centering
\includegraphics[width=8cm]{./figures/2cc53693047d69c0.png}
\caption{肉眼所见的天蝎座星座（已绘制星座连线）。} \label{fig_Scorpi_1}
\end{figure}
天蝎座包含许多明亮的恒星，其中包括心宿二 Antares（$\alpha$ Sco），意为“火星的对手”，因其显著的红色光泽而得名；$\beta_1$ Sco（Graffias 或 Acrab），一套三星系统；$\delta$ Sco（Dschubba，“额头”）；$\theta$ Sco（Sargas，名称源自苏美尔文化 $^{[3]}$）；$\nu$ Sco（Jabbah）；$\xi$ Sco；$\pi$ Sco（Fang）；$\sigma$ Sco（Alniyat）；以及 $\tau$ Sco（Paikauhale）。

位于蝎子弯曲尾部末端的是$\lambda$Sco（Shaula）与$\upsilon$Sco（Lesath），两者的名称均意为“毒刺”。由于它们彼此距离很近，$\lambda$Sco与$\upsilon$Sco 有时被合称为“猫眼” $^{[4]}$。

天蝎座中这些明亮恒星组成的形状类似于码头工人所用的钩子。其中大多数恒星都是最近的 OB 星协——天蝎—半人马星协（Scorpius--Centaurus）的巨大质量成员$^{[5]}$。

恒星$\delta$Sco 在长期保持稳定的 2.3 等亮度之后，于 2000 年 7 月在数周内突然增亮至 1.9 等。此后它成为一颗变星，亮度在 2.0 等与 1.6 等之间波动$^{[6]}$。这意味着在最亮时，它是天蝎座中第二亮的恒星。
\begin{figure}[ht]
\centering
\includegraphics[width=8cm]{./figures/a6e325eb5a989396.png}
\caption{按距离（红--绿三维视图）显示的天蝎座恒星分布，以及各恒星亮度（以恒星大小表示）} \label{fig_Scorpi_2}
\end{figure}
天蝎座 U 星（U~Scorpii）是目前已知爆发周期最短的新星，其周期约为 10 年$^{[7]}$。

天蝎座 AH 星是一颗红超巨星，也是已知体积最大的恒星之一，其半径约为太阳的 1400 倍。同时它也是一颗高光度恒星，亮度约为太阳的 340000 倍$^{[8]}$，但由于其视星等在 6.5 至 9.6 之间变化$^{[9]}$，肉眼通常无法直接观测到。

由$\omega^1$天蝎座 与$\omega^2$天蝎座 组成的一对近邻恒星是一对视双星，可用肉眼分辨。其中一颗是黄色巨星$^{[10]}$，另一颗则是天蝎—半人马星协中的一颗 B 型蓝色恒星$^{[11]}$。

曾被指定为$\gamma$天蝎座的恒星（尽管其实际位于天秤座边界之内）如今被称为$\sigma$天秤座。此外，在古希腊时期，整个天秤座曾被视为天蝎座的钳爪（Chelae~Scorpionis）；而在稍后的希腊时代，这些最西侧的恒星被重新诠释为由正义女神 Astraea（由相邻的处女座所代表）高举的一副天平。天秤座作为独立星座的划分在古希腊或古罗马时期被正式确立$^{[12]}$。
\subsubsection{深空天体}
\begin{figure}[ht]
\centering
\includegraphics[width=6cm]{./figures/36a84b2fd26c01cc.png}
\caption{天蝎座与银河系，同画面中可见心宿二（Antares）附近的 M4 与 M80，画面中央下方的 M6 与 M7，画面顶部的 NGC\,6124，以及画面中央正上方的 NGC\,6334。} \label{fig_Scorpi_3}
\end{figure}
由于其位置横跨银河系，这个星座包含了大量深空天体，例如疏散星团梅西耶~6（蝴蝶星团）与梅西耶~7（托勒密星团），位于 $\zeta^2$ 天蝎座附近的 NGC~6231，以及球状星团梅西耶~4 与梅西耶~80。

梅西耶~80（NGC~6093）是一座视星等为 7.3 的球状星团，距离地球约 33000 光年。它属于 Shapley~II 级球状星团；这一分类表明其核心高度集中且密度极高。M80 由查尔斯·梅西耶于$^{[1781]}$年发现。1860 年，阿图尔·冯·奥沃斯在此发现了新星天蝎座~T 星，这是一次相当罕见的发现$^{[13]}$。

NGC~6302，又称“甲虫星云”，是一座双极行星状星云。NGC~6334，又称“猫爪星云”，是一片发射星云及恒星形成区域。
\begin{figure}[ht]
\centering
\includegraphics[width=6cm]{./figures/5de286b7a5f3002e.png}
\caption{天蝎座的核心区域。画面中央偏左可见 M4。$\rho$蛇夫座分子云复合体的部分区域被心宿二（Antares）及其周围的恒星所照亮。} \label{fig_Scorpi_4}
\end{figure}
\subsection{神话}
\begin{figure}[ht]
\centering
\includegraphics[width=6cm]{./figures/adbb7e7cf91cc44d.png}
\caption{} \label{fig_Scorpi_5}
\end{figure}
在希腊神话中，与天蝎座相关的若干神话都将其归因于猎户座。根据其中一种说法，猎户座向女神阿尔忒弥斯及其母亲勒托夸口，说他将杀死地球上的所有动物。阿尔忒弥斯和勒托于是派出一只蝎子去杀死猎户座$^{[14]}$。他们的战斗引起了宙斯的注意，宙斯将双方都升上天空，以此提醒凡人要抑制过度的傲慢。在另一种版本的神话中，是阿尔忒弥斯的孪生兄弟阿波罗派出蝎子杀死猎户座，因为猎人因承认女神比自己更强而赢得了她的青睐。宙斯将猎户座和蝎子升上天空之后，前者每到冬季便在天空中追逐，而每到夏季，当蝎子出现时便逃离。在这两种版本中，都是阿尔忒弥斯请求宙斯将猎户座升上天空。

在另一则不涉及猎户座的希腊神话中，天上的蝎子曾在法厄同驾驶其父赫利俄斯的太阳战车时与其相遇$^{[15]}$。
\subsubsection{起源}  
巴比伦人称这一星座为 MUL.GIR.TAB——“蝎子”；该名称可以直译为“具有燃烧之刺的生物”$^{[16]}$。

在一些古老的记述中，天秤座被视为天蝎的钳爪。天秤座在巴比伦文献中被称为蝎之钳（zibānītu，参见阿拉伯语 zubānā），在希腊文中则称为 χηλαι$^{[17]}$。
\subsubsection{占星学}  
主条目：Scorpio（占星学）  
西方占星学中的天蝎座与天文学上的天蝎座星座并不相同。从天文学角度看，太阳仅在 11 月 23 日至 11 月 28 日这 6 天内位于国际天文学联合会定义的天蝎座边界之内。差异的很大一部分源于蛇夫座的存在，而该星座仅被少数占星家采用。天蝎座在印度占星体系中对应的宿为 Anuradha、Jyeshtha 和 Mula$^{[citation\ needed]}$。
